\chapter{Attention and Transformers}

\section{Attention Mechanisms}

Sequence models based on recurrence compress past information into a fixed-size hidden state. This creates a bottleneck when modeling long sequences.

\medskip

\noindent
\textbf{Core idea of attention.}  
Instead of relying on a single hidden state, attention allows the model to directly access all relevant parts of the input.

\medskip

\noindent
\textbf{Query-key-value framework.}  
Attention is defined in terms of three components:
\begin{itemize}
    \item Query ($q$)
    \item Keys ($k_i$)
    \item Values ($v_i$)
\end{itemize}

\medskip

\noindent
\textbf{Attention weights.}
\[
\alpha_i = \frac{\exp(\langle q, k_i \rangle)}{\sum_j \exp(\langle q, k_j \rangle)}.
\]

\medskip

\noindent
\textbf{Output.}
\[
\mathrm{Attn}(q, K, V) = \sum_i \alpha_i v_i.
\]

\medskip

\noindent
\textbf{Interpretation.}
\begin{itemize}
    \item The query determines which inputs are relevant
    \item Weights $\alpha_i$ represent importance
\end{itemize}

\medskip

\noindent
\textbf{Insight.}  
Attention replaces fixed memory with a flexible, data-dependent mechanism.

\section{Self-Attention}

Self-attention applies the attention mechanism within a single sequence.

\medskip

\noindent
\textbf{Definition.}  
Each element of the sequence attends to all other elements.

\medskip

\noindent
\textbf{Construction.}  
Given inputs $x_1, \dots, x_T$, define:
\[
q_i = W_Q x_i, \quad k_i = W_K x_i, \quad v_i = W_V x_i.
\]

\medskip

\noindent
\textbf{Attention output.}
\[
y_i = \sum_{j=1}^T \alpha_{ij} v_j,
\quad
\alpha_{ij} = \frac{\exp(\langle q_i, k_j \rangle)}{\sum_{l} \exp(\langle q_i, k_l \rangle)}.
\]

\medskip

\noindent
\textbf{Properties.}
\begin{itemize}
    \item Captures dependencies between all positions
    \item Fully parallelizable across sequence elements
\end{itemize}

\medskip

\noindent
\textbf{Positional encoding.}  
Since self-attention is permutation-invariant, positional information must be added explicitly.

\medskip

\noindent
\textbf{Insight.}  
Self-attention provides a flexible mechanism for modeling relationships within sequences.

\section{Transformer Architecture}

The transformer architecture is built entirely from attention mechanisms.

\medskip

\noindent
\textbf{Basic components.}
\begin{itemize}
    \item Multi-head self-attention
    \item Feedforward layers
    \item Residual connections
    \item Normalization layers
\end{itemize}

\medskip

\noindent
\textbf{Multi-head attention.}  
Multiple attention mechanisms are applied in parallel:
\[
\mathrm{MultiHead}(Q,K,V) = \mathrm{Concat}(H_1, \dots, H_h) W_O,
\]
where each head $H_i$ computes attention with different projections.

\medskip

\noindent
\textbf{Feedforward network.}
\[
\mathrm{FFN}(x) = \sigma(W_1 x + b_1) W_2 + b_2.
\]

\medskip

\noindent
\textbf{Residual structure.}
\[
x \leftarrow x + \mathrm{Layer}(x),
\]
followed by normalization.

\medskip

\noindent
\textbf{Architecture layout.}
\begin{itemize}
    \item Stack of identical blocks
    \item Each block: attention + feedforward
\end{itemize}

\medskip

\noindent
\textbf{Encoder--decoder structure.}
\begin{itemize}
    \item Encoder processes input sequence
    \item Decoder generates output sequence with attention to encoder outputs
\end{itemize}

\medskip

\noindent
\textbf{Insight.}  
Transformers replace recurrence with attention, enabling efficient modeling of long-range dependencies.

\section{Scaling and Applications}

Transformers have demonstrated remarkable scalability and performance across domains.

\medskip

\noindent
\textbf{Scaling behavior.}
\begin{itemize}
    \item Performance improves with model size, data, and computation
    \item Enables training of very large models
\end{itemize}

\medskip

\noindent
\textbf{Parallelization.}
\begin{itemize}
    \item Unlike RNNs, transformers process sequences in parallel
    \item Efficient use of modern hardware
\end{itemize}

\medskip

\noindent
\textbf{Applications.}
\begin{itemize}
    \item Natural language processing (translation, text generation)
    \item Vision (vision transformers)
    \item Speech processing
    \item Multimodal learning
\end{itemize}

\medskip

\noindent
\textbf{Pretraining and fine-tuning.}
\begin{itemize}
    \item Large models trained on massive datasets
    \item Adapted to specific tasks through fine-tuning
\end{itemize}

\medskip

\noindent
\textbf{Limitations.}
\begin{itemize}
    \item Quadratic complexity in sequence length
    \item High computational cost
\end{itemize}

\medskip

\noindent
\textbf{Insight.}  
Scaling and data availability are key drivers of modern deep learning performance.

\section{A Unifying Perspective}

Attention and transformers unify several ideas:

\begin{itemize}
    \item \textbf{Representation:} dynamic weighting of inputs
    \item \textbf{Computation:} parallel processing of sequences
    \item \textbf{Architecture:} modular and scalable design
\end{itemize}

\medskip

\noindent
\textbf{Core components.}
\begin{itemize}
    \item Query-key-value attention
    \item Multi-head mechanisms
    \item Residual and normalization layers
\end{itemize}

\medskip

\noindent
\textbf{Key insight.}  
Attention provides a flexible mechanism for modeling interactions, replacing fixed sequential structures.

\section{Outlook}

This chapter introduced attention mechanisms and transformers:
\begin{itemize}
    \item attention as a flexible memory mechanism,
    \item self-attention for sequence modeling,
    \item transformer architectures,
    \item scaling behavior and applications.
\end{itemize}

\medskip

In the next chapter, we study generative models, beginning with probabilistic modeling and latent variable approaches.

\medskip

\noindent
\textbf{Guiding principle.}  
Modern architectures succeed by enabling flexible, scalable interactions between data elements.